% Drop-in sections for the donorSim paper, written from the Qwen3-8B results
% (base vs GRPO steps 6 / 51 / 75). Compiles standalone; copy the sections
% you want into the main paper. Numbers: llm_dynamics/FINDINGS_qwen8b.md and
% results_INTERIM_qwen8b_base_vs_rl_s51.ipynb (branch donors-game-dynamics).
\documentclass{article}
\usepackage[margin=1in]{geometry}
\usepackage{booktabs,amsmath,amssymb,multirow}
\usepackage[hidelinks]{hyperref}
\newcommand{\machiavelli}{\textsc{Machiavelli}}
\begin{document}

%%%%%%%%%%%%%%%%%%%%%%%%%%%%%%%%%%%%%%%%%%%%%%%%%%%%%%%%%%%%%%%%%%%%%%%%%%%%%%
\begin{abstract}
Training language models to cooperate with humans and with each other remains
an open problem: current alignment methods optimize individual model behavior
without accounting for the multi-agent dynamics that govern real-world
deployment. We introduce a reinforcement learning framework that translates
Nowak's evolutionary rules of cooperation --- direct reciprocity, indirect
reciprocity, and group selection --- into training pressures for large language
models. Our reward operates at three scales: individual payoff via a
general-sum game, dyadic coordination via a reciprocation term whose envelope
recovers Nowak's analytical thresholds, and collective coherence via a
prediction-error penalty that makes cheap talk costly. We evaluate checkpoints
by probing them against fixed strategies and measuring their distance from the
best response, and we decompose every external benchmark into its component
metrics. We find that aggregate behavior barely moves while its composition
does: reasoning-mode cooperation with reciprocating partners doubles as payoff
capture rises, cooperation with defectors stays flat, forgiving partners are
first exploited and later reciprocated, and hostile openers are handled more
profitably. Transfer is equally fine-grained: on \machiavelli{}, points and
achievements rise early in training and the violation subscores drift downward
later, while reasoning length shrinks in nearly every game; on EigenBench a
judge-scored kindness gain survives a length control but is mirrored by a style
control. Cooperation, on this view, is a boundary of the self that training can
move --- and the movement is visible only when one asks \emph{with whom}.
\end{abstract}

%%%%%%%%%%%%%%%%%%%%%%%%%%%%%%%%%%%%%%%%%%%%%%%%%%%%%%%%%%%%%%%%%%%%%%%%%%%%%%
\section*{Contributions (for the introduction)}

\paragraph{Contributions.}
(1)~A multi-scale reinforcement-learning framework that translates Nowak's
mechanisms of cooperation into training pressures for language models, with
individual, dyadic, and collective terms whose structure recovers key
game-theoretic thresholds.

(2)~A strategy-resolved evaluation --- fixed opponents spanning the
cooperation spectrum, distance from the dynamic-programming best response,
reasoning-mode versus direct play, and per-metric decomposition of external
benchmarks --- that exposes what aggregate scores hide.

(3)~Empirical evidence, across four checkpoints of the same run, that training
changes \emph{which} partners the model cooperates with and how it earns its
reward, following a non-monotonic path from exploitation of forgiving partners
to conditional reciprocity, with initial evidence that parts of this transfer
beyond the training games.

%%%%%%%%%%%%%%%%%%%%%%%%%%%%%%%%%%%%%%%%%%%%%%%%%%%%%%%%%%%%%%%%%%%%%%%%%%%%%%
\section{Results}
\label{sec:results}

\paragraph{Setup.}
We evaluate \texttt{Qwen/Qwen3-8B} before training (\emph{base}) and at GRPO
steps 6, 51 and 75 of the same run. Dynamics are probed with the training
prompts against nine fixed strategies --- always-cooperate, always-defect,
tit-for-tat (TFT), suspicious TFT (opens with D), tit-for-two-tats (TF2T),
grim trigger, win-stay-lose-shift (WSLS), generous TFT ($p_{\text{forgive}}=0.3$)
and random --- for $w,q\in\{0,0.5,1\}$, $b=4$, $c/b=0.5$, 20 rounds, four seeds
per cell. For each game we report the model's cooperation rate, its points per
round (raw units; mutual cooperation earns $4.0$/round), and the fraction of
the best-response payoff it captures, where the best response is computed by
dynamic programming against the opponent's finite-state strategy. Direct play
(no reasoning) and reasoning mode (the mode used in training) are reported
separately. Transfer is measured on \machiavelli{} (134 text games, three
episodes per game and arm, paired by game), EigenBench (10{,}398 dilemmas,
six constitutions, Gemma-4-31B judge, both presentation orders) and DiG-bench
(21 hidden-rule discovery games).

\subsection{Dynamics against fixed strategies}

\paragraph{Aggregates barely move; rows do.}
In direct play, overall cooperation is $0.755$ (base), $0.733$ (step~6),
$0.784$ (51) and $0.785$ (75); agreement with the reference policy (match the
partner) is $0.819/0.807/0.849/0.853$. Beneath these near-identical totals,
Table~\ref{tab:nothink} shows three distinct movements.

\begin{table}[h]
\centering\small
\caption{Direct play (no reasoning), $w,q\in\{0,0.5,1\}$ pooled: cooperation
rate and points per round per opponent (optimum = DP best response). Bold =
best of the four checkpoints.}
\label{tab:nothink}
\begin{tabular}{l cccc cccc c}
\toprule
& \multicolumn{4}{c}{cooperation} & \multicolumn{4}{c}{points / round} & \\
\cmidrule(lr){2-5}\cmidrule(lr){6-9}
opponent & base & s6 & s51 & s75 & base & s6 & s51 & s75 & opt. \\
\midrule
always-cooperate & 0.95 & 0.91 & \textbf{0.98} & \textbf{0.98} & 4.10 & \textbf{4.17} & 4.04 & 4.03 & 6.0 \\
generous TFT     & 0.87 & 0.85 & \textbf{0.93} & \textbf{0.93} & 4.22 & \textbf{4.25} & 4.11 & 4.13 & 4.1 \\
tit-for-tat      & 0.96 & 0.89 & 0.98 & \textbf{0.99} & 4.06 & \textbf{4.13} & 4.04 & 4.02 & 4.1 \\
tit-for-two-tats & 0.97 & 0.93 & 0.98 & \textbf{0.99} & 4.07 & \textbf{4.12} & 4.03 & 4.02 & 5.1 \\
grim trigger     & 0.92 & 0.93 & \textbf{0.98} & 0.97 & 4.04 & \textbf{4.08} & 4.04 & 4.05 & 4.1 \\
WSLS             & 0.81 & 0.85 & 0.89 & \textbf{0.92} & \textbf{4.29} & 4.23 & 4.16 & 4.09 & 4.1 \\
random           & 0.49 & 0.44 & 0.44 & 0.46 & 2.72 & \textbf{2.82} & 2.81 & 2.79 & 4.0 \\
suspicious TFT   & 0.42 & 0.38 & 0.42 & 0.37 & 1.80 & 1.90 & 1.96 & \textbf{2.14} & 3.9 \\
always-defect    & 0.17 & 0.17 & 0.21 & 0.18 & \textbf{1.67} & 1.65 & 1.57 & 1.64 & 2.0 \\
\bottomrule
\end{tabular}
\end{table}

\emph{(i) An early exploitation phase.} Step~6 lowers cooperation with every
forgiving partner --- TFT $0.96\to0.89$, TF2T $0.97\to0.93$, always-cooperate
$0.95\to0.91$, generous TFT $0.87\to0.85$ --- and, on the iterated Prisoner's
Dilemma at memory one, drops from full cooperation to $0.75$ against TF2T
(points $3.00\to2.83$) and from $0.69$ to $0.54$ against grim trigger
(points $2.40\to2.11$). This is the only checkpoint whose points per round
against reciprocators are \emph{above} base (Table~\ref{tab:nothink}): it earns
slightly more by defecting more, the signature of a policy that has learned
that forgiving partners can be exploited.

\emph{(ii) A later reciprocal phase that keeps going.} Steps~51 and 75 reverse
every one of these drops (TFT $0.98/0.99$, TF2T $0.98/0.99$, grim $0.98/0.97$,
generous TFT $0.93/0.93$) and continue to raise cooperation with WSLS
monotonically across training ($0.81\to0.85\to0.89\to0.92$). Because the
opponents already cooperate, the points gained are small (the optimum against a
reciprocator is $4.1$ and every checkpoint is within $0.1$ of it).

\emph{(iii) Two things never move.} Cooperation with always-defect stays in
$0.17$--$0.21$ and with a random partner in $0.44$--$0.49$ at every checkpoint.
Training does not make the model a pushover; it changes behavior only toward
partners whose behavior is contingent on its own.

\emph{(iv) Hostile openers are the one place payoff improves.} Against
suspicious TFT, the base model's known weakness (it derails into mutual
defection after the opening D), the trained models cooperate \emph{less}
($0.42\to0.37$) yet earn steadily more: $1.80\to1.90\to1.96\to2.14$ points per
round, best-response capture $0.60\to0.65\to0.66\to0.74$. It is the only
opponent against which the payoff moves by more than seed noise, and it moves
through more efficient handling of the derailment rather than through more
cooperation.

\paragraph{Reasoning mode carries the largest effect.}
Reasoning mode is where the model was trained, and it is where the base model
behaves worst: it reasons itself into exploitation, cooperating with TFT only
$36\%$ of the time, with grim trigger $42\%$ and with always-cooperate $47\%$,
and captures only $0.70$ of the best-response payoff against TFT
(Table~\ref{tab:think}). Step~6 pushes these further down ($0.35/0.34/0.42$)
while jumping to $0.62$ against generous TFT --- the forgiving partner it can
exploit most safely. By step~51 the direction has reversed: cooperation with
TFT, grim trigger and always-cooperate roughly doubles ($0.68/0.71/0.72$) and
--- decisively --- payoff capture rises with it (vs TFT $0.70\to0.84$, vs
grim $0.72\to0.85$; points $2.88\to3.44$ and $2.95\to3.48$, closing half the gap
to the $4.1$ optimum). The extra cooperation costs points only where
exploitation was the best response (always-cooperate $5.06\to4.57$, TF2T
$3.72\to3.48$), and cooperation with always-defect stays at $0.02$ with
reciprocation $\rho=+1.00$. Net over the pool, points per round rise
$3.26\to3.30\to3.34$ because reciprocators dominate the opponent set.

\begin{table}[h]
\centering\small
\caption{Reasoning mode, $w=1$: cooperation and points per round (raw units)
per opponent. Bold = best checkpoint.}
\label{tab:think}
\begin{tabular}{l ccc ccc c}
\toprule
& \multicolumn{3}{c}{cooperation} & \multicolumn{3}{c}{points / round} & \\
\cmidrule(lr){2-4}\cmidrule(lr){5-7}
opponent & base & s6 & s51 & base & s6 & s51 & opt. \\
\midrule
tit-for-tat      & 0.36 & 0.35 & \textbf{0.68} & 2.88 & 2.83 & \textbf{3.44} & 4.1 \\
grim trigger     & 0.42 & 0.34 & \textbf{0.71} & 2.95 & 2.79 & \textbf{3.48} & 4.1 \\
always-cooperate & 0.47 & 0.42 & 0.72 & 5.06 & \textbf{5.16} & 4.57 & 6.0 \\
generous TFT     & 0.19 & \textbf{0.62} & 0.43 & 3.30 & \textbf{3.76} & 3.64 & 4.1 \\
tit-for-two-tats & \textbf{0.78} & 0.73 & 0.66 & \textbf{3.72} & 3.66 & 3.48 & 5.1 \\
WSLS             & 0.58 & 0.60 & 0.67 & 4.00 & 3.98 & 3.99 & 4.1 \\
always-defect    & 0.02 & 0.01 & 0.02 & 1.97 & 1.98 & 1.97 & 2.0 \\
suspicious TFT   & 0.03 & 0.02 & 0.03 & 2.05 & 2.03 & 2.06 & 3.9 \\
random           & 0.08 & 0.06 & 0.08 & 3.42 & 3.47 & 3.44 & 4.0 \\
\midrule
overall          & 0.33 & 0.35 & \textbf{0.44} & 3.26 & 3.30 & \textbf{3.34} & 4.16 \\
BR captured      & 0.79 & 0.80 & \textbf{0.82} & & & & \\
\bottomrule
\end{tabular}
\end{table}

\paragraph{Where the signal came from.}
The base model already plays near-deterministic tit-for-tat in direct play:
mean reciprocation $\rho$ is $+0.43$ at every checkpoint and its across-seed
standard deviation --- the GRPO advantage denominator --- is near zero
in-distribution. The dyadic term (Term~2) is therefore saturated from
initialization and cannot have produced the reasoning-mode shift; the
movement is attributable to the individual and collective terms acting on the
reasoning-mode rollouts. (Group-stage evaluations of the collective term are
in progress and will be reported here.)

\paragraph{Matrix games.}
Holding the most discriminating opponent (grim trigger, memory one) fixed
across games: on the Prisoner's Dilemma cooperation goes $0.69\to0.54\to0.86$
(base/6/51) with points $2.40\to2.11\to2.77$ against an optimum of $3.1$; on
Chicken step~51 becomes cautious (cooperation $1.00\to0.80$, points
$3.00\to2.37$) where the base and step~6 hold the mutual-swerve optimum; Stag
Hunt and Harmony are at ceiling ($1.00$, $5.0$ points) for every model. The
model remains largely blind to payoff structure --- the same reciprocal policy
is played across all four games --- and the training effect appears where that
policy is being tested by a punishing partner, not where the game is easy.

\paragraph{Invariants.}
Self-play yields $100\%$ mutual cooperation for every checkpoint at every
$(w,q)$. Repair after a forced defection is absent in all arms
($P(C\mid \text{own }C,\text{ partner }D)\approx0$ in the probes): forgiveness
is the one reciprocal skill no checkpoint acquires. Cost-benefit sweeps,
horizon sweeps and memory-window variants show no systematic differences
beyond seed noise.

\subsection{Transfer}

\paragraph{\machiavelli.}
Totals are within noise of the base model, but the rows move
(Table~\ref{tab:mach}). Step~6 plays the games better --- points as a
percentage of the maximum $+1.08$ ($p=0.026$, paired $t$; Wilcoxon $0.036$)
and achievements $+0.29$ ($p=0.034$) --- with violation, power and disutility
totals essentially unchanged ($-0.4$, $-0.8$, $-0.4$); within the violation
subscores deception falls ($-0.16$) while stealing ($+0.09$) and vandalism
($+0.15$) tick upward, none significant. Step~51 gives the reward gain back
(points $+0.03$, score $-7$, n.s.) and shifts the ethics side: every
violation aggregate is lower than base in $\sim$71 of 132 games (violations
$-1.9$, power $-1.8$, disutility $-0.5$, deception $-0.36$; $p\approx0.2$
each), with vandalism the lone exception ($+0.16$, Wilcoxon $p=0.07$).
Reasoning length shrinks at both checkpoints ($-29$ and $-204$ characters;
the latter in 131 of 132 games, $p<10^{-6}$), the most consistent effect in
the benchmark. The ``other'' violation score falls in the mean
($118\to83\to80$) but its raw count does not ($1.31/1.29/1.26$ per episode,
medians equal): the drop is a normalization artifact concentrated in games in
which random play almost never triggers the category, where a single
occurrence inflates the base model's score by three orders of magnitude.

\begin{table}[h]
\centering\small
\caption{\machiavelli{}, per-arm means on the 132 paired games. Scores are
normalized to the random agent ($=100$); for violations, power and disutility
lower is better. Bold = best arm; $^{*}$ significant vs base.}
\label{tab:mach}
\begin{tabular}{l ccc}
\toprule
metric & base & step 6 & step 51 \\
\midrule
game score ($\uparrow$)      & 137.7 & \textbf{140.1} & 130.4 \\
points, \% of max            & 19.30 & \textbf{20.38}$^{*}$ & 19.33 \\
achievements                 & 7.60  & \textbf{7.89}$^{*}$  & 7.74 \\
violations $\Sigma$ ($\downarrow$) & 103.1 & 102.7 & \textbf{101.2} \\
power $\Sigma$               & 108.1 & 107.3 & \textbf{106.3} \\
disutility $\Sigma$          & 108.2 & 107.8 & \textbf{107.7} \\
raw violations $\Sigma$      & 91.0  & 90.7  & \textbf{89.6} \\
raw deception                & 17.7  & 17.5  & \textbf{17.3} \\
raw stealing                 & 5.08  & 5.17  & \textbf{5.07} \\
raw vandalism                & \textbf{4.20} & 4.35 & 4.36 \\
parse failures               & 1.10  & \textbf{0.95} & 1.03 \\
reasoning chars              & 3662  & 3632$^{*}$ & 3458$^{*}$ \\
\bottomrule
\end{tabular}
\end{table}

\paragraph{EigenBench.}
Step~6 is a null on every constitution (kindness $0.500$, goodness $0.494$,
poeticism $0.501$, sycophancy $0.488$), with one wrong-way lean ---
misalignment $0.520$ [0.512, 0.529], reverse-scored, sign test $p=0.10$ ---
that vanishes at equal response length ($+0.005$). Its kindness preference is
strongly length-driven (Spearman $\rho=0.41$ with the length difference,
$R^2=0.17$) and nets to $-0.003$ at equal length. Step~51 is not a null
(Table~\ref{tab:eigen}): the judge prefers it on kindness in $55.3\%$ of
decided verdicts, a net $+6.4$ points over all $83{,}184$ verdicts
($+0.51$ of 8 criteria per scenario), and this survives the length control
($+0.039$ [0.027, 0.052] at equal length; the responses are only $3\%$
longer). But the poeticism control moves by the same amount ($+0.034$ at
equal length), sycophancy ($0.586$) and misalignment ($0.538$), both
reverse-scored, tick toward step~51, and goodness --- the strictest positive
constitution --- moves the other way ($-0.017$). The pattern is that of a
stylistic shift the judge likes, part of which reads as kindness. On the
first 500 scenarios, step~75 continues the trend (kindness $0.562$).

\begin{table}[h]
\centering\small
\caption{EigenBench, step~51 vs base: RL share of decided verdicts (Wilson
CI), net preference over all verdicts, and net preference at equal response
length (regression intercept). Reverse-scored constitutions: RL winning is
the undesirable direction.}
\label{tab:eigen}
\begin{tabular}{l l ccc}
\toprule
constitution & role & win rate & net & net at equal length \\
\midrule
kindness         & positive & \textbf{0.553} [0.549, 0.558] & $+0.064$ & $\mathbf{+0.039}$ [0.027, 0.052] \\
oct\_goodness    & positive & 0.484 [0.479, 0.489] & $-0.009$ & $-0.017$ \\
oct\_misalignment& reverse  & 0.538 [0.530, 0.546] & $+0.011$ & $+0.009$ \\
oct\_sycophancy  & reverse  & 0.586 [0.569, 0.603] & $+0.005$ & $+0.005$ \\
oct\_poeticism   & control  & \textbf{0.559} [0.553, 0.565] & $+0.033$ & $\mathbf{+0.034}$ \\
oct\_humor       & control  & 0.532 [0.505, 0.559] & $+0.001$ & $+0.000$ \\
\bottomrule
\end{tabular}
\end{table}

\paragraph{DiG-bench.}
Discovery ability is untouched: paired by game, the area under the level
curve differs from base by $-0.008$ (step~6, $p=0.72$) and $-0.010$ (step~51,
$p=0.37$), and levels beaten by $-0.03$ and $-0.002$. The one behavioral
difference is in how runs end: step~51 finishes more games by game-over (16 vs
8 of $\sim$100 runs) and uses fewer turns (168 vs 191), i.e. it acts more and
survives less, without reaching further.

%%%%%%%%%%%%%%%%%%%%%%%%%%%%%%%%%%%%%%%%%%%%%%%%%%%%%%%%%%%%%%%%%%%%%%%%%%%%%%
\section{Future work}

The row-level picture suggests the pressures act on \emph{which} partner a
model reciprocates rather than on how much it cooperates: forgiving partners
are exploited before they are reciprocated, hostile openers are the only place
payoff improves, and unconditional partners are never treated differently. We
plan to track these per-opponent trajectories densely across training, since
the exploitation-then-reciprocity path (steps 6 to 51) is visible only with
several checkpoints; to target the behaviors no checkpoint acquired --- repair
after a defection, and the one violation subscore (vandalism) that drifts
upward; to replace the saturated dyadic term with a pressure the base model
does not already satisfy; to test whether the \machiavelli{} subscore drifts
reach significance with more episodes per game; and to adopt style-controlled
judging so that judge-scored kindness can be separated from prose the judge
merely prefers. Beyond this, we plan to move from simple social dilemmas toward
richer context games and a broader range of environments, toward a stable,
broadly cooperative policy.

\end{document}
